\subsection{TikZ -- Block Diagrams}
\label{subsec:tikz-control-systems}
For control-system block diagrams, TikZ requires:

\begin{lstlisting}[style=latex]
\usepackage{tikz}
\usetikzlibrary{positioning, calc}
\end{lstlisting}


% ============================================================
\subsubsection{Basic Block Diagram}

\begin{lstlisting}[style=latex]
\begin{tikzpicture}
    \node[input] (input) {$R(s)$};
    \node[block, right=of input] (G1) {$s+4$};
    \node[block, right=of G1] (G2)
        {$\dfrac{10}{s^2+16s+20}$};
    \node[output, right=of G2] (output) {$Y(s)$};

    \draw[->] (input) -- (G1);
    \draw[->] (G1) -- (G2);
    \draw[->] (G2) -- (output);
\end{tikzpicture}
\end{lstlisting}

\begin{center}
\begin{tikzpicture}
    \node[input] (input) {$R(s)$};
    \node[block, right=of input] (G1) {$s+4$};
    \node[block, right=of G1] (G2)
        {$\dfrac{10}{s^2+16s+20}$};
    \node[output, right=of G2] (output) {$Y(s)$};

    \draw[->] (input) -- (G1);
    \draw[->] (G1) -- (G2);
    \draw[->] (G2) -- (output);
\end{tikzpicture}
\end{center}


% ============================================================
\subsubsection{Node Positioning}

\begin{lstlisting}[style=latex]
right=of node
left=of node
above=of node
below=of node

right=1cm of node
below=2\dnd cm of node % 2 times the default node distance
\end{lstlisting}

The first form uses the default \texttt{node distance}; the second
overrides it locally.


% ============================================================
\subsubsection{Node Types}

\begin{lstlisting}[style=latex]
\node[input]   (input) {$R(s)$};

\node[block]   (G) {$G(s)$};

\node[sum]     (sum1) {};

\node[pickoff] (p1) {};

\node[output]  (output) {$Y(s)$};
\end{lstlisting}


% ============================================================
\subsubsection{Connections}

\begin{lstlisting}[style=latex]
% Straight line
\draw[->] (A) -- (B);

% Straight line without arrow
\draw[-] (A) -- (B);

% Vertical, then horizontal
\draw[->] (A) |- (B);

% Horizontal, then vertical
\draw[->] (A) -| (B);
\end{lstlisting}

Node anchors can be selected explicitly:

\begin{lstlisting}[style=latex]
node.north
node.south
node.east
node.west
\end{lstlisting}


% ============================================================
\subsubsection{Summing Junctions}
Like a connection, but with a plus or minus sign at the junction.
The following code snippet will be used after positioning the nodes:
\begin{lstlisting}[style=latex]
\node[sum] (sum1) {};

\draw[->] (G)
    -- node[plus, above left] {}
    (sum1);

\draw[->] (H)
    -| node[minus, below left] {}
    (sum1.south);
\end{lstlisting}

\begin{center}
\begin{tikzpicture}
    \node[input] (input) {$R(s)$};
    \node[sum, right=of input] (sum1) {};
    \node[block, right=of sum1] (G) {$G(s)$};

    \coordinate (mid) at ($(sum1)!0.5!(G)$);
    \node[block, below=of mid] (H) {$H(s)$};

    \draw[->] (input)
        -- node[plus, above left] {}
        (sum1);

    \draw[->] (sum1) -- (G);

    \draw[->] (G.south) |- (H.east);

    \draw[->] (H.west)
        -| node[minus, below left] {}
        (sum1.south);
\end{tikzpicture}
\end{center}


% ============================================================
\subsubsection{Pickoff Points}

\begin{lstlisting}[style=latex]
\node[pickoff, right=of G] (p1) {};

\draw[-]  (G) -- (p1);
\draw[->] (p1) -- (output);
\draw[->] (p1.south) |- (feedback.east);
\end{lstlisting}


% ============================================================
\subsubsection{Feedback Loops}

\begin{lstlisting}[style=latex]
\coordinate (mid)
    at ($(sum1)!0.5!(pickoff1)$);

\node[block, below=of mid] (feedback)
    {$H(s)$};

\draw[->] (pickoff1.south)
    |- (feedback.east);

\draw[->] (feedback.west)
    -| node[minus, below left] {}
    (sum1.south);
\end{lstlisting}

\begin{center}
\begin{tikzpicture}
    \node[input] (input) {$R(s)$};
    \node[sum, right=of input] (sum1) {};
    \node[block, right=of sum1] (plant) {$G(s)$};
    \node[pickoff, right=of plant] (pickoff1) {};
    \node[output, right=of pickoff1] (output) {$Y(s)$};

    \coordinate (mid)
        at ($(sum1)!0.5!(pickoff1)$);

    \node[block, below=of mid] (feedback)
        {$H(s)$};

    \draw[->] (input)
        -- node[plus, above left] {}
        (sum1);

    \draw[->] (sum1) -- (plant);
    \draw[-]  (plant) -- (pickoff1);
    \draw[->] (pickoff1) -- (output);

    \draw[->] (pickoff1.south)
        |- (feedback.east);

    \draw[->] (feedback.west)
        -| node[minus, below left] {}
        (sum1.south);
\end{tikzpicture}
\end{center}


% ============================================================
\subsubsection{Useful Coordinates}

Point halfway between two nodes:

\begin{lstlisting}[style=latex]
\coordinate (mid)
    at ($(A)!0.5!(B)$);
\end{lstlisting}

Other positions:

\begin{lstlisting}[style=latex]
($(A)!0.25!(B)$)   % 25 %
($(A)!0.50!(B)$)   % 50 %
($(A)!0.75!(B)$)   % 75 %
\end{lstlisting}
Where A and B are the names of the nodes. The \texttt{calc} library must be loaded for this to work.

% ============================================================
\subsubsection{Scaling}
In case the diagram doesn't fit in the page, we can scale it down using the \texttt{scale} option in the \texttt{tikzpicture} environment:
\begin{lstlisting}[style=latex]
\begin{tikzpicture}[scale=0.8, transform shape]
    ...
\end{tikzpicture}
\end{lstlisting}